\documentclass[12pt,a4paper]{article}
\usepackage[utf8]{inputenc}
\usepackage[T1]{fontenc}
\usepackage[french]{babel}
\usepackage[margin=2.5cm]{geometry}
\usepackage{graphicx}
\usepackage{listings}
\usepackage{xcolor}
\usepackage{hyperref}
\usepackage{titlesec}
\usepackage{amsmath}
\usepackage{amssymb}
\usepackage{enumitem}
\usepackage{fancyhdr}
\usepackage{booktabs}
\usepackage{array}
\usepackage{longtable}
\usepackage{multirow}
\usepackage{caption}
\usepackage{subcaption}

% Configuration des couleurs pour les listings
\definecolor{codegreen}{rgb}{0,0.6,0}
\definecolor{codegray}{rgb}{0.5,0.5,0.5}
\definecolor{codepurple}{rgb}{0.58,0,0.82}
\definecolor{backcolour}{rgb}{0.95,0.95,0.92}

\lstdefinestyle{mystyle}{
    backgroundcolor=\color{backcolour},   
    commentstyle=\color{codegreen},
    keywordstyle=\color{magenta},
    numberstyle=\tiny\color{codegray},
    stringstyle=\color{codepurple},
    basicstyle=\ttfamily\footnotesize,
    breakatwhitespace=false,         
    breaklines=true,                 
    captionpos=b,                    
    keepspaces=true,                 
    numbers=left,                    
    numbersep=5pt,                  
    showspaces=false,                
    showstringspaces=false,
    showtabs=false,                  
    tabsize=2,
    frame=single,
    framesep=10pt
}

\lstset{style=mystyle}

% Configuration des en-têtes et pieds de page
\pagestyle{fancy}
\fancyhf{}
\fancyhead[L]{Projet : Contrôle d'Incubateur à Œufs IoT}
\fancyhead[R]{\thepage}
\renewcommand{\headrulewidth}{0.5pt}

% Configuration des titres
\titleformat{\section}[block]{\normalfont\Large\bfseries}{\thesection}{1em}{}
\titleformat{\subsection}[block]{\normalfont\large\bfseries}{\thesubsection}{1em}{}
\titleformat{\subsubsection}[block]{\normalfont\normalsize\bfseries}{\thesubsubsection}{1em}{}

% Informations du rapport
\title{
    \vspace{2cm}
    \textbf{Rapport de Projet : Système de Contrôle d'Incubateur à Œufs IoT} \\
    \large Utilisation de l'ESP32, MQTT et Node.js pour la surveillance et le contrôle à distance
}
\author{Ghassen Saidi}
\date{\today}

\begin{document}

% Page de titre
\begin{titlepage}
    \centering
    \vspace*{1cm}
    
    \includegraphics[width=0.3\textwidth]{logo.png} % Remplacez par votre logo
    
    \vspace{1.5cm}
    
    {\huge\bfseries Rapport de Projet \\[0.5cm]
    Système de Contrôle d'Incubateur à Œufs IoT}
    
    \vspace{1.5cm}
    
    {\Large\bfseries Ghassen Saidi}
    
    \vspace{1cm}
    
    {\large\bfseries Projet de Fin d'Études\\[0.3cm]
    Département de Génie Électrique et Informatique}
    
    \vspace{1.5cm}
    
    \begin{flushleft}
        \textbf{Encadrant :} Pr. [Nom de l'Encadrant] \\
        \textbf{Date :} \today \\
        \textbf{Version :} 1.0
    \end{flushleft}
    
    \vfill
    
    {\large Année Universitaire 2023-2024}
    
\end{titlepage}

% Table des matières
\tableofcontents
\newpage

% Liste des figures
\listoffigures
\newpage

% Liste des tableaux
\listoftables
\newpage

% Introduction
\section{Introduction et Objectifs}

\subsection{Contexte et Problématique}
L'incubation artificielle des œufs est une pratique essentielle en aviculture moderne, permettant de contrôler précisément les conditions environnementales nécessaires au développement embryonnaire. Cependant, les incubateurs traditionnels présentent plusieurs limitations : manque de surveillance en temps réel, impossibilité de contrôle à distance, et absence d'alertes automatiques en cas de défaillance.

Ce projet vise à résoudre ces problématiques en développant un système de contrôle intelligent basé sur les technologies IoT (Internet des Objets), offrant une surveillance continue et un contrôle précis des paramètres critiques pour l'incubation.

\subsection{Objectifs du Projet}
Les objectifs principaux de ce projet sont :

\begin{enumerate}[label=\textbf{\arabic*.}]
    \item \textbf{Surveillance en temps réel} : Mesurer et afficher en continu la température et l'humidité à l'intérieur de l'incubateur.
    \item \textbf{Contrôle automatique} : Maintenir automatiquement les conditions optimales (37.5°C, 60\% humidité) pour l'incubation.
    \item \textbf{Interface de contrôle à distance} : Développer une interface web accessible depuis n'importe quel appareil connecté à Internet.
    \item \textbf{Système multi-tâches} : Gérer simultanément plusieurs capteurs et actionneurs avec une architecture modulaire.
    \item \textbf{Communication robuste} : Implémenter un système de communication fiable entre l'ESP32 et le serveur web.
\end{enumerate}

\subsection{Technologies Utilisées}
Le projet utilise les technologies suivantes :

\begin{table}[h]
    \centering
    \begin{tabular}{|l|l|p{8cm}|}
        \hline
        \textbf{Technologie} & \textbf{Version} & \textbf{Rôle dans le projet} \\
        \hline
        ESP32 & DevKit C V4 & Microcontrôleur principal avec WiFi intégré \\
        \hline
        MQTT & 3.1.1 & Protocole de communication léger pour IoT \\
        \hline
        Node.js & 18.x & Plateforme serveur pour l'interface web \\
        \hline
        Socket.IO & 4.8.3 & Communication bidirectionnelle en temps réel \\
        \hline
        HTML/CSS/JavaScript & ES6 & Interface utilisateur web responsive \\
        \hline
        Wokwi & 1.0 & Environnement de simulation pour circuits électroniques \\
        \hline
        ArduinoJson & 6.x & Sérialisation/désérialisation JSON sur ESP32 \\
        \hline
        Express.js & 4.22.1 & Framework web pour Node.js \\
        \hline
    \end{tabular}
    \caption{Technologies utilisées dans le projet}
    \label{tab:technologies}
\end{table}

\newpage

% Description du circuit
\section{Description du Circuit}

\subsection{Schéma de Connexions}
Le circuit a été conçu et simulé dans l'environnement Wokwi. La configuration JSON complète est présentée ci-dessous :

\begin{lstlisting}[language=json, caption=Configuration JSON du circuit Wokwi]
{
  "version": 1,
  "author": "Ghassen Saidi",
  "editor": "wokwi",
  "parts": [
    { "type": "board-esp32-devkit-c-v4", "id": "esp", "top": 0, "left": 0, "attrs": {} },
    { "type": "wokwi-led", "id": "led1", "top": 150, "left": -63.4, "attrs": { "color": "red" } },
    { "type": "wokwi-servo", "id": "servo1", "top": 132.4, "left": 230.4, "attrs": {} },
    {
      "type": "wokwi-led",
      "id": "led2",
      "top": 66,
      "left": 150.6,
      "rotate": 90,
      "attrs": { "color": "blue" }
    },
    {
      "type": "wokwi-buzzer",
      "id": "bz1",
      "top": 21.6,
      "left": 241.8,
      "attrs": { "volume": "0.1" }
    },
    { "type": "wokwi-potentiometer", "id": "pot1", "top": -10.9, "left": -105.8, "attrs": {} },
    { "type": "wokwi-dht22", "id": "dht1", "top": 19.5, "left": 407.4, "attrs": {} }
  ],
  "connections": [
    [ "esp:TX", "$serialMonitor:RX", "", [] ],
    [ "esp:RX", "$serialMonitor:TX", "", [] ],
    [ "dht1:SDA", "esp:4", "green", [ "v9.6", "h-67.1", "v-19.2" ] ],
    [ "esp:5V", "dht1:VCC", "red", [ "h-33.41", "v28.8", "h268.8" ] ],
    [ "dht1:GND", "esp:GND.1", "black", [ "v86.4", "h-326.4", "v48", "h-201.6", "v-115.2" ] ],
    [ "pot1:VCC", "esp:5V", "red", [ "v0" ] ],
    [ "pot1:GND", "esp:GND.1", "black", [ "v0" ] ],
    [ "pot1:SIG", "esp:VP", "green", [ "v19.2", "h47.6", "v-28.8" ] ],
    [ "led1:C", "esp:GND.1", "black", [ "v48", "h-28.4", "v-86.4" ] ],
    [ "servo1:PWM", "esp:5", "green", [ "h-19.2", "v-76.6" ] ],
    [
      "servo1:GND",
      "esp:GND.1",
      "black",
      [ "h-105.6", "v48", "h0", "v38.4", "h-201.6", "v-28.8" ]
    ],
    [ "servo1:V+", "esp:5V", "red", [ "h-67.2", "v38.5", "h-192", "v-28.8" ] ],
    [ "esp:18", "led2:A", "green", [ "h0" ] ],
    [ "led2:C", "esp:GND.3", "black", [ "h0" ] ],
    [ "bz1:2", "esp:19", "green", [ "v19.2", "h-58", "v-48", "h-86.4", "v19.2" ] ],
    [ "bz1:1", "esp:GND.3", "black", [ "v9.6", "h-38.4", "v-48", "h-105.6", "v9.6" ] ],
    [ "esp:2", "led1:A", "green", [ "h14.44", "v86.4", "h-134.4", "v-57.6" ] ]
  ],
  "dependencies": {}
}
\end{lstlisting}

\subsection{Description des Composants}

\subsubsection{Capteurs}
\begin{table}[h]
    \centering
    \begin{tabular}{|p{4cm}|p{3cm}|p{8cm}|}
        \hline
        \textbf{Composant} & \textbf{Spécifications} & \textbf{Fonction} \\
        \hline
        DHT22 & Précision : ±0.5°C, ±2-5\% & Mesure la température et l'humidité intérieure \\
        \hline
        Potentiomètre & 10kΩ, linéaire & Permet un contrôle manuel de la température de consigne \\
        \hline
    \end{tabular}
    \caption{Spécifications des capteurs utilisés}
    \label{tab:capteurs}
\end{table}

\subsubsection{Actionneurs}
\begin{table}[h]
    \centering
    \begin{tabular}{|p{4cm}|p{3cm}|p{8cm}|}
        \hline
        \textbf{Composant} & \textbf{Spécifications} & \textbf{Fonction} \\
        \hline
        Servomoteur SG90 & 0°-180°, 2.5kg/cm & Contrôle l'ouverture de la fenêtre de ventilation \\
        \hline
        LED bleue (ventilateur) & 5mm, 20mA & Simule le fonctionnement du ventilateur de refroidissement \\
        \hline
        Buzzer passif & 5V, 2kHz & Émet une alarme en cas d'humidité excessive \\
        \hline
        LED rouge (état) & 5mm, 20mA & Indicateur visuel de l'état du système \\
        \hline
    \end{tabular}
    \caption{Spécifications des actionneurs utilisés}
    \label{tab:actionneurs}
\end{table}

\subsection{Alimentation et Masse}
Le système utilise une alimentation unique de 5V fournie par l'ESP32. Tous les composants partagent la même masse (GND) pour éviter les problèmes de masse flottante. Des résistances de limitation de courant (220Ω) sont intégrées pour protéger les LEDs.

\begin{figure}[h]
    \centering
    \includegraphics[width=0.8\textwidth]{circuit_schema.png} % Remplacez par votre schéma
    \caption{Schéma électrique du circuit}
    \label{fig:circuit}
\end{figure}

\newpage

% Explications du code
\section{Explications du Code et Choix Techniques}

\subsection{Architecture du Système}

\subsubsection{Structure Multi-Topics MQTT}
L'architecture MQTT utilise plusieurs topics spécialisés pour optimiser la communication :

\begin{lstlisting}[language=c++, caption=Définition des topics MQTT dans Firmware.ino]
// Topics MQTT définis dans Firmware.ino
const char* topic_heartbeat = "incubator/heartbeat";
const char* topic_sensor_readings = "incubator/sensor_readings";
const char* topic_control_setpoints = "incubator/control_setpoints";
const char* topic_manual_controls = "incubator/manual_controls";
const char* topic_system_status = "incubator/system_status";
const char* topic_commands = "incubator/commands";
\end{lstlisting}

\textbf{Avantages de cette approche :}
\begin{itemize}
    \item \textbf{Séparation des préoccupations} : Chaque topic gère un type de données spécifique
    \item \textbf{Optimisation du trafic réseau} : Réduction des données redondantes
    \item \textbf{Évolutivité} : Facilité d'ajouter de nouveaux capteurs/actionneurs
    \item \textbf{Débogage facilité} : Monitoring spécifique par topic
\end{itemize}

\subsubsection{Communication Bidirectionnelle}
Le système implémente une communication bidirectionnelle complète :
\begin{itemize}
    \item \textbf{ESP32 → Serveur} : Données de capteurs, états système, heartbeat
    \item \textbf{Serveur → ESP32} : Commandes de contrôle, requêtes de données, paramètres
\end{itemize}

\subsection{Code ESP32 (Firmware.ino)}

\subsubsection{Initialisation et Configuration}
\begin{lstlisting}[language=c++, caption=Fonction setup() du firmware ESP32]
void setup() {
  // Initialisation série pour le débogage
  Serial.begin(115200);
  
  // Configuration des broches GPIO
  pinMode(LED_PIN, OUTPUT);
  pinMode(FAN_PIN, OUTPUT);
  pinMode(BUZZER_PIN, OUTPUT);
  pinMode(POT_PIN, INPUT);
  
  // Initialisation du capteur DHT22
  dht.begin();
  
  // Configuration du servomoteur
  windowServo.setPeriodHertz(50);
  windowServo.attach(SERVO_PIN, 500, 2400);
  windowServo.write(0);
  
  // Connexion WiFi avec gestion d'erreurs
  setup_wifi();
  
  // Configuration MQTT avec callback
  client.setServer(mqtt_server, mqtt_port);
  client.setCallback(callback);
}
\end{lstlisting}

\subsubsection{Boucle Principale avec Timing Précis}
\begin{lstlisting}[language=c++, caption=Fonction loop() avec gestion du timing]
void loop() {
  // Reconnexion MQTT si nécessaire
  if (!client.connected()) {
    reconnect();
  }
  client.loop();
  
  unsigned long now = millis();
  
  // Lecture des capteurs toutes les 2 secondes
  if (now - lastMsg > 2000) {
    lastMsg = now;
    readSensors();
    readPotentiometer();
    runAutoControl();
    sendAllData();
  }
  
  // Heartbeat toutes les 5 secondes
  if (now - lastHeartbeat > 5000) {
    lastHeartbeat = now;
    sendHeartbeat();
  }
  
  delay(10); // Prévention du watchdog
}
\end{lstlisting}

\subsubsection{Contrôle Automatique avec Hystérésis}
\begin{lstlisting}[language=c++, caption=Algorithme de contrôle automatique]
void runAutoControl() {
  if (autoMode) {
    // Contrôle de température avec hystérésis de 3°C
    if (temperature >= maxTemperature) {
      if (!fanState) {
        digitalWrite(FAN_PIN, HIGH);
        fanState = true;
        setServoAngle(90);  // Ouvre la fenêtre
        windowState = true;
        digitalWrite(LED_PIN, HIGH);
        Serial.println("🔄 AUTO: Temperature high! Cooling activated.");
      }
    } else if (temperature <= (maxTemperature - 3)) {
      if (fanState) {
        digitalWrite(FAN_PIN, LOW);
        fanState = false;
        setServoAngle(0);  // Ferme la fenêtre
        windowState = false;
        digitalWrite(LED_PIN, LOW);
        Serial.println("🔄 AUTO: Temperature normal. Cooling deactivated.");
      }
    }
  }
}
\end{lstlisting}

\subsection{Serveur Node.js (index.js)}

\subsubsection{Architecture Serveur}
\begin{lstlisting}[language=javascript, caption=Configuration du serveur Node.js]
// Configuration Express + Socket.IO + MQTT
const express = require('express');
const mqtt = require('mqtt');
const http = require('http');
const socketIo = require('socket.io');

// Initialisation Express
const app = express();
const server = http.createServer(app);
const io = socketIo(server, {
  cors: {
    origin: "*",
    methods: ["GET", "POST"]
  }
});

// Middleware
app.use(cors());
app.use(express.json());
app.use(express.static('public'));

// Stockage des données en mémoire
let sensorData = {
  temperature: 25.0,
  humidity: 60.0,
  potentiometer: 0,
  servoPosition: 0,
  // ... autres données
};

// Connexion MQTT
const mqttClient = mqtt.connect(MQTT_BROKER, {
  clientId: 'nodejs_server_' + Math.random().toString(16).substr(2, 8),
  clean: true,
  connectTimeout: 4000,
  reconnectPeriod: 1000
});
\end{lstlisting}

\subsubsection{Gestion des Connexions WebSocket}
\begin{lstlisting}[language=javascript, caption=Gestion des événements Socket.IO]
io.on('connection', (socket) => {
  const clientId = socket.id.substring(0, 6);
  console.log(`🔌 New client connected: ${clientId}`);
  
  // Envoi des données initiales
  socket.emit('full-update', sensorData);
  
  // Gestion des commandes
  socket.on('control-command', (command) => {
    console.log(`🎮 Command from ${clientId}:`, command);
    
    // Publication sur MQTT
    mqttClient.publish(TOPICS.COMMANDS, JSON.stringify(command), 
      { qos: 1 }, (err) => {
        if (err) {
          socket.emit('command-error', { 
            message: 'Failed to send command',
            command: command 
          });
        } else {
          socket.emit('command-success', { 
            message: 'Command sent successfully',
            command: command 
          });
        }
    });
  });
});
\end{lstlisting}

\subsection{Interface Web (index.html)}

\subsubsection{Design Responsive}
L'interface utilise CSS Grid et Flexbox pour un design adaptatif :

\begin{lstlisting}[language=html, caption=Structure HTML responsive]
<div class="dashboard">
  <!-- Sensor Readings Card -->
  <div class="card">
    <div class="card-header">
      <h2 class="card-title">📊 Sensor Readings</h2>
    </div>
    <div class="sensor-grid">
      <div class="sensor-item">
        <div class="sensor-label">Temperature</div>
        <div class="sensor-value" id="tempValue">--</div>
        <div class="sensor-unit">°C</div>
      </div>
      <!-- ... autres capteurs -->
    </div>
  </div>
  <!-- ... autres cartes -->
</div>
\end{lstlisting}

\subsubsection{Gestion d'État Client}
\begin{lstlisting}[language=javascript, caption=Gestion d'état dans le client JavaScript]
// Stockage local des données
let currentData = {
  temperature: 25.0,
  humidity: 60.0,
  servoPosition: 0,
  maxTemperature: 30,
  maxHumidity: 80,
  mode: "AUTO",
  // ... autres états
};

// Mise à jour synchronisée des interfaces
function updateServoPositionDisplay(servoPosition) {
  // Mise à jour dans la carte des lectures
  document.getElementById('servoValue').textContent = servoPosition;
  
  // Mise à jour dans la carte d'état système
  document.getElementById('systemServo').textContent = servoPosition + '°';
  
  // Mise à jour du marqueur visuel
  updateServoMarker(servoPosition);
  
  // Mise à jour du statut texte
  updateServoStatus(servoPosition);
}
\end{lstlisting}

\subsection{Choix Techniques Justifiés}

\subsubsection{MQTT vs HTTP/REST}
\begin{table}[h]
    \centering
    \begin{tabular}{|p{6cm}|p{6cm}|}
        \hline
        \textbf{MQTT (choisi)} & \textbf{HTTP/REST (non choisi)} \\
        \hline
        Faible consommation de bande passante & Charge plus importante par requête \\
        \hline
        Communication publish/subscribe légère & Requêtes request/response plus lourdes \\
        \hline
        Support natif des connexions persistantes & Nécessite des techniques de polling ou websockets \\
        \hline
        Meilleur pour les dispositifs IoT à ressources limitées & Plus adapté aux applications web traditionnelles \\
        \hline
    \end{tabular}
    \caption{Comparaison MQTT vs HTTP/REST}
    \label{tab:mqtt_vs_http}
\end{table}

\subsubsection{Architecture Multi-Topics}
\begin{itemize}
    \item \textbf{\texttt{sensor\_readings}} : Données brutes des capteurs (température, humidité, potentiomètre, servo)
    \item \textbf{\texttt{control\_setpoints}} : Paramètres de configuration (température max, humidité max, mode)
    \item \textbf{\texttt{manual\_controls}} : États des actionneurs (ventilateur, fenêtre, LED, buzzer)
    \item \textbf{\texttt{system\_status}} : Santé du système (état capteur, RSSI WiFi, mémoire)
    \item \textbf{\texttt{heartbeat}} : Connectivité et informations basiques
    \item \textbf{\texttt{commands}} : Commandes du serveur vers l'ESP32
\end{itemize}

\subsubsection{Gestion des Erreurs et Robustesse}
\begin{itemize}
    \item \textbf{Reconnexion automatique} : WiFi et MQTT avec backoff exponentiel
    \item \textbf{Valeurs simulées} : En cas d'erreur du capteur DHT22
    \item \textbf{Timeout de connexion} : Détection de déconnexion ESP32 après 10 secondes
    \item \textbf{Validation JSON} : Parsing sécurisé avec gestion d'erreurs
    \item \textbf{QoS MQTT} : Niveau 1 pour les commandes importantes
\end{itemize}

\newpage

% Améliorations IoT
\section{Améliorations et Évolutions IoT}

\subsection{Sécurité et Accès à Distance avec Ngrok}

\subsubsection{Problématique}
Actuellement, l'interface web est accessible uniquement en réseau local. Pour un accès à distance sécurisé, nous proposons l'intégration de \textbf{Ngrok}.

\subsubsection{Solution Proposée : Ngrok avec HTTPS}
\begin{lstlisting}[language=bash, caption=Configuration de Ngrok]
# Installation de Ngrok
npm install -g ngrok

# Authentification avec votre token
ngrok authtoken <votre_token_ngrok>

# Tunnel HTTPS vers le serveur local
ngrok http 3000 --host-header="localhost:3000"

# Configuration avec domaine personnalisé
ngrok http 3000 --hostname=votre-domaine.ngrok.io
\end{lstlisting}

\subsubsection{Configuration ESP32 pour HTTPS}
\begin{lstlisting}[language=c++, caption=Support TLS/SSL sur ESP32]
// Configuration pour MQTT sécurisé
const char* mqtt_server = "broker.emqx.io";
const int mqtt_port = 8883; // Port TLS/SSL

// Certificat racine pour EMQX
const char* root_ca = \
"-----BEGIN CERTIFICATE-----\n" \
"MIIDrzCCApegAwIBAgIQCDvgVpBCRrGhdWrJWZHHSjANBgkqhkiG9w0BAQUFADBh\n" \
"... (certificat complet) ...\n" \
"-----END CERTIFICATE-----\n";

// Client WiFi sécurisé
WiFiClientSecure espClient;
PubSubClient client(espClient);

void setup() {
  // Configuration du certificat
  espClient.setCACert(root_ca);
  
  // Configuration des identifiants MQTT
  const char* mqtt_username = "user";
  const char* mqtt_password = "secure_password";
  
  // Connexion avec authentification
  client.connect(client_id, mqtt_username, mqtt_password);
}
\end{lstlisting}

\subsection{Autres Améliorations Possibles}

\subsubsection{Sécurité Renforcée avec JWT}
\begin{lstlisting}[language=javascript, caption=Authentification JWT côté serveur]
const jwt = require('jsonwebtoken');

// Middleware d'authentification
function authenticateToken(req, res, next) {
  const authHeader = req.headers['authorization'];
  const token = authHeader && authHeader.split(' ')[1];
  
  if (!token) return res.sendStatus(401);
  
  jwt.verify(token, process.env.JWT_SECRET, (err, user) => {
    if (err) return res.sendStatus(403);
    req.user = user;
    next();
  });
}

// Protection des routes API
app.get('/api/status', authenticateToken, (req, res) => {
  res.json({ status: 'authenticated', data: sensorData });
});
\end{lstlisting}

\subsubsection{Base de Données pour l'Historique}
\begin{lstlisting}[language=javascript, caption=Intégration MongoDB]
const mongoose = require('mongoose');

// Modèle de données
const sensorReadingSchema = new mongoose.Schema({
  temperature: { type: Number, required: true },
  humidity: { type: Number, required: true },
  timestamp: { type: Date, default: Date.now },
  deviceId: { type: String, required: true }
});

// Sauvegarde périodique
setInterval(() => {
  const reading = new SensorReading({
    temperature: sensorData.temperature,
    humidity: sensorData.humidity,
    deviceId: sensorData.client_id
  });
  
  reading.save()
    .then(() => console.log('Reading saved to database'))
    .catch(err => console.error('Error saving reading:', err));
}, 300000); // Toutes les 5 minutes
\end{lstlisting}

\subsubsection{Notifications Push}
\begin{lstlisting}[language=javascript, caption=Service de notifications]
const webpush = require('web-push');

// Configuration VAPID
webpush.setVapidDetails(
  'mailto:admin@incubator.com',
  process.env.VAPID_PUBLIC_KEY,
  process.env.VAPID_PRIVATE_KEY
);

// Stockage des abonnements
let pushSubscriptions = [];

// Envoi de notification d'alerte
function sendTemperatureAlert(currentTemp, threshold) {
  const message = `Alerte : Température ${currentTemp}°C (seuil : ${threshold}°C)`;
  
  pushSubscriptions.forEach(subscription => {
    webpush.sendNotification(subscription, JSON.stringify({
      title: '⚠️ Alerte Incubateur',
      body: message,
      icon: '/assets/icon-192.png',
      badge: '/assets/badge-72.png'
    })).catch(error => {
      console.error('Push notification error:', error);
    });
  });
}
\end{lstlisting}

\subsubsection{Progressive Web App (PWA)}
\begin{lstlisting}[language=json, caption=Manifest PWA]
{
  "name": "Incubator Controller",
  "short_name": "Incubator",
  "description": "Contrôle d'incubateur à œufs IoT",
  "start_url": "/",
  "display": "standalone",
  "background_color": "#16213e",
  "theme_color": "#4cc9f0",
  "icons": [
    {
      "src": "/icon-192.png",
      "sizes": "192x192",
      "type": "image/png"
    },
    {
      "src": "/icon-512.png",
      "sizes": "512x512",
      "type": "image/png"
    }
  ]
}
\end{lstlisting}

\subsection{Évolutions Matérielles}

\begin{table}[h]
    \centering
    \begin{tabular}{|p{5cm}|p{10cm}|}
        \hline
        \textbf{Composant} & \textbf{Description} \\
        \hline
        Capteur CO₂ SCD30 & Mesure la concentration de CO₂ pour la qualité de l'air \\
        \hline
        Caméra ESP32-CAM & Surveillance visuelle des œufs et détection d'éclosion \\
        \hline
        Module RTC DS3231 & Horloge temps réel pour la programmation précise \\
        \hline
        Module relais 4 canaux & Contrôle d'appareils 220V (humidificateur, chauffage) \\
        \hline
        Batterie Li-Po 18650 & Alimentation de secours avec gestion de charge \\
        \hline
        Écran OLED 0.96" & Affichage local des paramètres sans interface web \\
        \hline
    \end{tabular}
    \caption{Évolutions matérielles possibles}
    \label{tab:evolutions_materielles}
\end{table}

\subsection{Diagramme d'Architecture Amélioré}

\begin{figure}[h]
    \centering
    \includegraphics[width=0.9\textwidth]{architecture_amelioree.png} % Remplacez par votre diagramme
    \caption{Architecture système améliorée avec Ngrok et services cloud}
    \label{fig:architecture_amelioree}
\end{figure}

\newpage

% Conclusion
\section{Conclusion}

\subsection{Résumé des Réalisations}
Ce projet a permis de développer avec succès un système complet de contrôle d'incubateur à œufs IoT, intégrant :

\begin{itemize}
    \item ✅ Une architecture matérielle modulaire avec ESP32
    \item ✅ Un système de communication robuste basé sur MQTT
    \item ✅ Une interface web temps réel avec Socket.IO
    \item ✅ Un contrôle automatique avec hystérésis
    \item ✅ Une simulation complète dans l'environnement Wokwi
    \item ✅ Une documentation technique complète
\end{itemize}

\subsection{Perspectives et Améliorations}
Les principales pistes d'amélioration identifiées sont :

\begin{enumerate}
    \item \textbf{Sécurité renforcée} : Intégration de Ngrok pour l'accès HTTPS sécurisé à distance
    \item \textbf{Stockage des données} : Base de données pour l'historique et l'analyse
    \item \textbf{Application mobile} : Développement d'une application React Native
    \item \textbf{Intelligence artificielle} : Algorithmes de prédiction et d'optimisation
    \item \textbf{Automatisation avancée} : Programmation des cycles d'incubation
\end{enumerate}

\subsection{Impact et Applications}
Ce système peut être adapté à diverses applications nécessitant un contrôle précis de l'environnement :

\begin{itemize}
    \item \textbf{Agriculture} : Serres intelligentes, élevage
    \item \textbf{Laboratoire} : Incubateurs bactériologiques, étuves
    \item \textbf{Domestique} : Cave à vin, rangement des aliments
    \item \textbf{Industriel} : Chambres climatiques, processus de séchage
\end{itemize}

\subsection{Remerciements}
Je tiens à remercier [Nom de l'encadrant] pour son accompagnement et ses précieux conseils tout au long de ce projet, ainsi que les enseignants du département pour leur transmission de connaissances.

\vspace{1cm}

\begin{center}
    \textbf{Fin du rapport}
\end{center}

% Annexes
\appendix
\section{Annexes}

\subsection{Code Source Complet}
Les codes sources complets sont disponibles sur GitHub :\\
\url{https://github.com/[votre-compte]/egg-incubator-iot}

\subsection{Références Bibliographiques}
\begin{enumerate}
    \item ESP32 Datasheet, Espressif Systems, 2023
    \item MQTT Version 3.1.1 Specification, OASIS, 2014
    \item Node.js Documentation, OpenJS Foundation, 2023
    \item Wokwi Documentation, Wokwi, 2023
    \item "IoT System Architecture: Design and Evaluation", Springer, 2022
\end{enumerate}

\subsection{Glossaire}
\begin{description}
    \item[IoT] Internet of Things (Internet des Objets)
    \item[MQTT] Message Queuing Telemetry Transport
    \item[API] Application Programming Interface
    \item[PWA] Progressive Web Application
    \item[JWT] JSON Web Token
    \item[REST] Representational State Transfer
    \item[SSL/TLS] Secure Sockets Layer / Transport Layer Security
    \item[GPIO] General Purpose Input/Output
    \item[PWM] Pulse Width Modulation
    \item[ADC] Analog to Digital Converter
\end{description}

\end{document}
